% ============================================================================
% PLANTILLA 01 — EXAMEN DE DESARROLLO (teórico)
% ----------------------------------------------------------------------------
% Para: exámenes parciales/finales/sustitutorios con preguntas de desarrollo,
% definiciones, conceptualización, comparaciones, respuesta corta y ensayo.
% Compilar: latexmk -pdf plantilla-01-examen-desarrollo.tex
% Solucionario: cambie la línea \documentclass por
%   \documentclass[soluciones]{academic-exam}
% ============================================================================
\documentclass{academic-exam}

% ------------------------------------------------------- metadatos ----------
\universidad{Universidad Nacional de San Cristóbal de Huamanga}
\facultad{Facultad de Ciencias Económicas, Administrativas y Contables}
\escuela{Escuela Profesional de Economía}
\curso{Macroeconomía II}
\codigocurso{EC-344}
\docente{Mg. Nombre Apellido}
\periodo{Semestre 2026-I}
\tipoevaluacion{Examen Parcial}
\fechaevaluacion{21 de julio de 2026}
\duracion{110 minutos}
\modalidad{Presencial}

\begin{document}
\membrete

\begin{instrucciones}
\begin{itemize}
  \item Lea cada pregunta con atención antes de responder. La claridad,
        el orden y la ortografía forman parte de la calificación.
  \item Responda únicamente en los espacios asignados, con lapicero azul
        o negro. No se aceptan respuestas a lápiz.
  \item Está prohibido el uso de apuntes, libros o dispositivos electrónicos.
  \item El examen consta de \totalpreguntas{} preguntas y suma
        \totalpuntos{} puntos.
\end{itemize}
\end{instrucciones}


% ============================================================ PREGUNTAS ====
\begin{pregunta}[4]
Defina con precisión los siguientes conceptos y proporcione un ejemplo
aplicado a la economía peruana para cada uno:
\begin{apartados}
  \item Política monetaria no convencional. \pts{2}
  \item Trampa de liquidez. \pts{2}
\end{apartados}
\lineasrespuesta{8}
\begin{solucion}
\textbf{a)} La política monetaria no convencional agrupa los instrumentos
que emplea el banco central cuando la tasa de interés de referencia se
aproxima a su límite inferior efectivo: expansión cuantitativa, forward
guidance y operaciones de reporte ampliadas. \emph{Ejemplo:} las
operaciones de reporte de cartera del BCRP durante 2020 (programa
Reactiva Perú).

\textbf{b)} La trampa de liquidez es la situación en la que la demanda de
dinero se vuelve perfectamente elástica a la tasa de interés vigente, de
modo que aumentos de la oferta monetaria no reducen la tasa ni estimulan
la demanda agregada.
\end{solucion}
\end{pregunta}

\begin{pregunta}[4]
Compare el enfoque keynesiano y el enfoque monetarista respecto de la
efectividad de la política fiscal. Organice su respuesta en un cuadro con,
al menos, tres criterios de comparación (supuestos, mecanismo de
transmisión y evidencia empírica).
\espaciorespuesta{7cm}
\begin{solucion}
Una respuesta completa debe contrastar: (i) \emph{supuestos} — rigideces
nominales y demanda efectiva frente a flexibilidad de precios y expectativas
adaptativas; (ii) \emph{mecanismo} — multiplicador fiscal mayor que uno
frente a efecto expulsión (\emph{crowding out}) vía tasa de interés;
(iii) \emph{evidencia} — multiplicadores altos en recesiones profundas
(Auerbach y Gorodnichenko, 2012) frente a multiplicadores bajos en
economías abiertas con tipo de cambio flexible.
\end{solucion}
\end{pregunta}

\begin{pregunta}[6]
\enquote{La dolarización financiera limita la potencia de la política
monetaria en economías emergentes.} Analice críticamente esta afirmación
considerando el caso peruano entre 2002 y 2024. Su respuesta debe incluir:
\begin{apartados}
  \item El canal de hoja de balance y el descalce cambiario. \pts{2}
  \item El papel del esquema de metas explícitas de inflación. \pts{2}
  \item Al menos una política de desdolarización aplicada por el BCRP. \pts{2}
\end{apartados}
\lineasrespuesta{12}
\begin{solucion}
El descalce cambiario hace que una depreciación deteriore la hoja de
balance de los deudores en dólares con ingresos en soles, lo que genera un
canal contractivo que compite con el canal expansivo tradicional del tipo
de cambio. El esquema de metas de inflación (desde 2002) ancló las
expectativas y redujo el traspaso cambiario, mientras que los encajes
diferenciados en moneda extranjera y el programa de repos de sustitución
(2014--2016) aceleraron la desdolarización del crédito, que pasó de
alrededor de 80\,\% a menos de 25\,\%.
\end{solucion}
\end{pregunta}

\begin{pregunta}[6]
Redacte un ensayo breve (máximo dos carillas) sobre el siguiente tema:
\emph{¿Debe el Banco Central de Reserva incorporar objetivos de empleo en
su mandato?} Tome una posición y defiéndala con argumentos teóricos y
evidencia empírica.

\begin{criterios}
\begin{tabularx}{\linewidth}{@{}X c@{}}
  \textbf{Criterio} & \textbf{Puntaje}\\
  \arrayrulecolor{grislinea}\hline\addlinespace[3pt]
  Tesis clara y posición definida            & 1{,}5 \\
  Argumentación teórica consistente          & 2{,}0 \\
  Uso pertinente de evidencia empírica       & 1{,}5 \\
  Redacción, estructura y ortografía         & 1{,}0 \\
\end{tabularx}
\end{criterios}
\lineasrespuesta{10}
\begin{solucion}
No existe una única respuesta correcta: se califica la solidez de la
argumentación según la rúbrica. Una respuesta de nivel excelente discute
el mandato dual de la Reserva Federal, la divina coincidencia
(Blanchard y Galí, 2007), los riesgos de dominancia fiscal y la evidencia
sobre credibilidad de bancos centrales con mandato único.
\end{solucion}
\end{pregunta}

\findelexamen
\end{document}
